\documentclass[10pt,a4paper]{amsart}
\usepackage[margin=19mm]{geometry}
\usepackage{amsmath,amssymb,mathtools}
\usepackage[scr=rsfs,cal=euler]{mathalfa}
\usepackage{xfrac}
\usepackage[unicode,hidelinks,bookmarks=false]{hyperref}
\newcommand{\ie}{i.e.\ }
\newcommand{\cf}{cf.\ }
\newcommand{\resp}{respectively\ }
\newcommand{\Subsection}{\subsection}
\providecommand{\nobreakdash}{\nobreak\hbox{-}}
\setlength{\emergencystretch}{2em}
\multlinegap0pt
\usepackage[shortlabels]{enumitem}
\usepackage{amssymb}
\usepackage{mathtools}
\usepackage{bm}
\newtheorem{theorem}{Theorem}
\newcommand{\R}{\mathbb R}
\title{\textbf{Endpoint Structure in Three-Dimensional Navier--Stokes}}
\author{Rishad Shahmurov}
\date{Source: arXiv:2606.07875v2 (CC BY 4.0)}
\begin{document}
\maketitle

\noindent\emph{Source and licence.} \textbf{Endpoint Structure in Three-Dimensional Navier--Stokes}. Version arXiv:2606.07875v2 (CC BY 4.0), distributed by arXiv under the Creative Commons Attribution 4.0 International licence, http://creativecommons.org/licenses/by/4.0/, as recorded in the arXiv licence metadata for this exact version and shown on the abstract page https://arxiv.org/abs/2606.07875v2. Full licence notice: \texttt{licenses/arxiv-2606.07875-CC-BY-4.0.txt}.\par\smallskip

\noindent\textit{Editorial scope.} Counted unit: the article's theorem theorem (source0.tex lines 320--340). The article's complete original proof of it is delivered with it (source0.tex lines 342--475, one \texttt{proof} environment of 134 lines). No mathematics is written, repaired or re-derived: the statement and the proof are the article's own lines, in the article's own notation, and no retained line is substituted or re-flowed.  No further source-local context is needed: every cross-reference of every retained line resolves inside the two delivered spans.  Nothing was omitted from any retained span, so the package carries no omission record.  No retained line contains a citation, so the wrapper carries no bibliography and no bibliography entry.  No retained line contains a figure environment, an image-inclusion command or drawing code, so no image is delivered and the package declares no asset dependency.  Editorial formatting only, in addition: an AMS \texttt{amsart} preamble that restates the article's own declarations and macros used by the retained lines, namely the environments \texttt{theorem} on separate counters, together with the standard packages those lines need.  The retained environments print in this document as Theorem 1 in order of appearance, whereas the article prints the same items with the numbering of its own full document.  The complete native source of the article as distributed in the arXiv e-print for this exact version is bundled unchanged as \texttt{sources/202286/source0.tex}.
\begin{theorem}
Let \((u_n,p_n^\sharp)\) be a normalized sequence of suitable weak
Navier--Stokes solutions on a fixed parabolic chamber.  Assume uniform
scale-invariant local energy and dissipation bounds, local
\(L^{3/2}\)-pressure bounds in the fixed root gauge, local boundedness of the
critical threshold measures, and tightness/boundedness of the explicitly
retained nonnegative defect coordinates.  Then after a subsequence the
sequence generates an endpoint package \(\mathcal E_{\rm NS}\).  Moreover:
\begin{enumerate}[label=(\alph*),leftmargin=*]
\item \(u_n\to U\) strongly in \(L^q_{\rm loc}\) for every
      \(q<10/3\), after the standard suitable-solution compactness step;
\item the localized near pressure converges strongly in \(L^r_{\rm loc}\) for every
      \(1<r<5/3\), while harmonic/tail parts are retained separately;
\item nonnegative defect measures are weak-* compact and lower semicontinuous;
\item weak nonlinear-product loss is represented by \(\mathcal R\);
\item direction oscillation and amplitude concentration are represented by
      \(\nu_\xi\) and \(\mu_{\rm conc}\);
\item the fixed pressure projection may be applied before passage to the limit,
      so \(\mathcal J_{\rm triad}\) is retained distributionally.
\end{enumerate}
\end{theorem}

\begin{proof}
We divide the extraction into the velocity, pressure, measure, and signed-current
coordinates.  All compactness statements below are first proved on a strict
subchamber \(Q'\Subset Q\) and are then diagonalized over a countable exhaustion
of the normalized chamber.

\emph{Step 1: energy bounds and the subcritical spacetime exponent.}
The assumed local energy and dissipation bounds give, uniformly in \(n\),
\[
 u_n\quad\hbox{bounded in}\quad
 L^\infty(I';L^2(B'))\cap L^2(I';H^1(B'))
\]
for every \(Q'=B'\times I'\Subset Q\).  The three-dimensional Sobolev
embedding \(H^1(B')\hookrightarrow L^6(B')\), combined with interpolation in
space and time, yields the standard estimate
\[
 \|u_n\|_{L^{10/3}(Q')}
 \le C_{Q'}
 \|u_n\|_{L^\infty_tL^2_x(Q')}^{2/5}
 \|u_n\|_{L^2_tH^1_x(Q')}^{3/5}.
\]
Thus \((u_n)\) is uniformly bounded at the natural suitable-weak exponent
\(10/3\).

\emph{Step 2: a time-derivative bound.}
Choose a cutoff \(\chi\in C_c^\infty(Q)\) that is one on \(Q'\).  On the support
of \(\chi\), the equation can be written schematically as
\[
 \partial_t u_n
 =\nu\Delta u_n-\mathbb P\nabla\!\cdot(u_n\otimes u_n),
\]
where the Leray projection is harmless in negative Sobolev spaces.  The first
term is bounded in \(L^2(I';H^{-1}(B'))\).  Since
\(u_n\in L^{10/3}(Q')\), one has
\(u_n\otimes u_n\in L^{5/3}(Q')\), and therefore the nonlinear term is bounded
in a space of the form
\(L^{5/3}(I';W^{-1,5/3}(B'))\).  After lowering exponents if necessary, both
terms are bounded in one reflexive negative Sobolev space
\[
 \partial_tu_n\quad\hbox{bounded in}\quad
 L^\sigma(I';W^{-1,m}(B'))
\]
for some fixed \(\sigma>1\) and \(m>1\).  This is the compactness input needed
for the Aubin--Lions lemma.

\emph{Step 3: strong local velocity convergence.}
Aubin--Lions applied to
\(H^1(B')\Subset L^2(B')\hookrightarrow W^{-1,m}(B')\) gives, after passage to
a subsequence,
\[
 u_n\to U\qquad\hbox{strongly in }L^2(Q').
\]
The uniform \(L^{10/3}\) bound and interpolation then imply
\[
 u_n\to U\qquad\hbox{strongly in }L^q(Q')
 \quad\text{for every }1\le q<10/3.
\]
In particular, if \(2<q<10/3\), then
\[
 u_n\otimes u_n\to U\otimes U
 \qquad\hbox{strongly in }L^{q/2}(Q').
\]
This observation is useful because it identifies precisely which quadratic
velocity products are already compact and which additional nonlinear
coordinates must be retained only for weaker limiting procedures.

\emph{Step 4: the near pressure.}
Fix a spatial cutoff \(\eta\) supported strictly inside the root chamber and
identically one on the spatial projection of \(Q'\).  The near pressure is
\[
 p_n^{\rm near}=\mathcal R_i\mathcal R_j(\eta u_{n,i}u_{n,j}).
\]
Choose any \(1<r<5/3\) and set \(q=2r\); then \(2<q<10/3\).  By Step 3 the
localized product converges strongly in \(L^r\).  The Riesz transforms are
bounded on \(L^r(\R^3)\) for \(1<r<\infty\), hence
\[
 p_n^{\rm near}\to
 \mathcal R_i\mathcal R_j(\eta U_iU_j)
 \qquad\hbox{strongly in }L^r(Q').
\]
The portion of \(p_n^\sharp\) not contained in this near field is harmonic on
the smaller spatial region, up to the explicitly retained tail and time-face
contributions.  Those pieces are therefore kept as separate endpoint
coordinates rather than being absorbed into the near-pressure convergence.
The use of the fixed root representative is essential here: no scale-dependent
additive pressure constant is introduced during extraction.

\emph{Step 5: nonnegative defect measures.}
Let \((\mu_n)\) denote any of the nonnegative Radon coordinates for which the
hypotheses give a uniform mass bound on compact subsets.  By the Riesz
representation theorem and Banach--Alaoglu, the unit ball of finite Radon
measures is weak-* compact on every compact subchamber.  Hence, after a
subsequence,
\[
 \mu_n\stackrel{*}{\rightharpoonup}\mu.
\]
For every nonnegative \(\varphi\in C_c\),
\[
 \int\varphi\,d\mu
 \le \liminf_{n\to\infty}\int\varphi\,d\mu_n
\]
whenever the coordinate arises from a lower-semicontinuous convex energy or
dissipation density.  This gives the required lower semicontinuity for the
local-energy, concentration, and critical-threshold coordinates.

\emph{Step 6: Reynolds/product defects and generalized oscillation data.}
Whenever a nonlinear quantity is already covered by the strong convergence in
Step 3, its defect is zero on \(Q'\).  For any additional bounded nonlinear
observable for which only weak convergence is available, we define its defect
as the difference between the weak limit of the observable and the observable
formed from the weak limit.  By assumption these defect coordinates are
uniformly bounded and tight, so they admit subsequential limits.  Directional
oscillations are encoded by the Young measure \(\nu_\xi\), and loss through
large amplitudes is encoded by \(\mu_{\rm conc}\).  The point of retaining both
objects is that an oscillatory defect and an amplitude-concentration defect
have different mathematical meanings and cannot be reconstructed from one
another after the limit has been taken.

\emph{Step 7: signed currents.}
The pressure-projected local-triad current is not inserted into a positive
measure budget.  Its test operator, cutoff, pressure representative, and
finite-dimensional quotient are fixed before the subsequence is selected.
Therefore, for every compactly supported test field \(\phi\), its pairing with
\(\phi\) is a bounded scalar sequence.  Diagonal extraction over a countable
dense family of test fields yields a distributional limit, denoted
\(\mathcal J_{\rm triad}\).  The same principle applies to any other retained
signed distributional coordinate.

Finally, perform a diagonal subsequence extraction over the countable
exhaustion \(Q'_1\Subset Q'_2\Subset\cdots\Subset Q\) and over the finite or
countable collection of endpoint coordinates.  Every component of
\(\mathcal E_{\rm NS}\) is then defined on one common subsequence, which is the
required faithful endpoint package.
\end{proof}

\end{document}
